% Verbatim copy of /Users/alingalatan/workspaces/grothe/handoff/manuscript/legacy/ch25-parts/92-joint-sixth-formal.tex (file dated 2026-10-25). Only this line was added.
\section{A common contraction budget for the sixth moment and Stein cuts}
\label{sec:joint-sixth-formal-20261022}
\paragraph{Scope.}
The following dimension-independent inequalities are now proved in Lean for genuine
complex $(2,1)$ Wick cubics. They are additional constraints on the \emph{same}
tensor used by the Stein cut. They do not assume conditional Gaussianity,
independent factors, or a bounded coordinate core. The finite-node port checks every majorant and every use of its kurtosis
range. The complete native Lean replay of the lifted cover passes, with
$B=0.7134218$ and $K_G^{\mathbb C}>1.401695322458$ under the existing
node-cap hypothesis.

Normalize $\E|q|^2=1$, put $\chi=\E|q|^4-2$ and $x=\sqrt\chi$, and use the
five contraction invariants $(n,h,b,y,z)$ of the contraction-Gram proof. Set
\[
 s_0=8n+2h+8b+8z,\qquad
 h_0=\frac{s_0+24(3y-z)}3.
\]
Both are nonnegative, despite the negative $z$ coefficient when $h_0$ is
expanded, and
\[
 \chi=\frac83s_0+4h_0.
\]
For every $\alpha>0$, the two proved bounds are
\begin{align}
 \E|q|^6&\le6+9\chi+x\{150\alpha s_0+(198+150/\alpha)h_0\},
 \label{eq:joint-sixth-formal-triangle}\\
 \E|q|^6&\le6+9\chi+x\{(135\alpha/2)s_0+
 (11781/32+135/(2\alpha))h_0\}.
 \label{eq:joint-sixth-formal-young}
\end{align}
The normalization bridge is explicit: if $a=\Re q$ and $f$ is its real
symmetric cubic kernel, then $\E a^2=1/2$ and
$s_0=1296\,\operatorname{tenS}(f)$,
$h_0=1296\,\operatorname{tenH}(f)$.
Using the variance-one real kernel $\sqrt2 f$ instead changes the factor to
$324$. Mixing these two conventions would invalidate the numerical cut.

\paragraph{Why these are joint constraints.}
Let $P,Q,V,C$ denote the four linear contraction expressions
\[
\begin{split}
 P&=16n+4h+16b/3+32y,\quad Q=64b/3+32y,\\
 V&=(80n+20h+16b+96y+32z)/9,\quad C=(128b+96y)/9.
\end{split}
\]
Here $V$ is the Stein variance and $C$ the squared complex fluctuation norm.
A numerical node previously certified $\lambda V+\mu C\le AP+BQ$ by two
positive semidefinite Gram multipliers and five nonnegative coefficient
remainders. For either positive bracket $J=A_\alpha s_0+B_\alpha h_0$ above,
reserve $\zeta x_*J$ in that \emph{same} coefficient identity:
\[
 \lambda V+\mu C+\zeta x_*J\le AP+BQ,\qquad \zeta\ge0.
\]
In coordinates $(n,h,b,y,z)$ the reserved vector is
\[
 \zeta x_*\bigl(8A_\alpha+8B_\alpha/3,
 2A_\alpha+2B_\alpha/3,
 8A_\alpha+8B_\alpha/3,
 24B_\alpha,
 8A_\alpha-16B_\alpha/3\bigr).
\]
Each remainder is checked \emph{after} subtracting this vector. Its last
entry may be negative; a componentwise nonnegative replacement would lose
valid information.

On a checked range $0\le x\le x_*$, the moment inequality gives
$\E|q|^6\le6+9\chi+x_*J$. Therefore the old radial cut $G(u)$ can be
replaced by
\[
 G(u)+\zeta(-u^3+9u^2-12).
\]
Indeed, its expectation is the old completion squares plus the two
nonnegative slacks
\[
 AP+BQ-\lambda V-\mu C-\zeta x_*J,
 \qquad
 \zeta(x_*J+6+9\chi-\E|q|^6).
\]
This is a reusable certificate interface: one tensor budget pays for
both terms before any scalar law relaxation. Adding two separately
optimized numerical savings would not establish this identity.

\paragraph{Formal proof locations and remaining computation.}
The Lean modules are \code{Grothe/LowerAnalytic/JointSixth.lean}
(\code{gram\_j1}, \code{gram\_j2}, \code{gram\_j3}, and the two normalization
bridges) and \code{Grothe/Lower/Reserve.lean}
(\code{reserved\_gram\_budget}, \code{joint\_reserve\_value},
\code{joint\_moment\_range}, \code{reserved\_moment\_cut}). The third module \code{Grothe/LowerAnalytic/JointBudget.lean}
(\code{joint\_sixth\_budget}) connects this finite acceptance to the actual four
spectral expressions of a circular cubic on the checked range. Their axiom
checks use only \code{propext}, \code{Classical.choice}, and \code{Quot.sound}.
The real tensor proof retains the negative cubic-trace term and uses
\[
 -20v+\frac{165}2\sqrt{vw}\le\frac{5445}{64}w
\]
by completing the square. Positivity of the symmetrized contraction
operator is not assumed.

The full node/range/cover integration is in the lower checker, with
\code{Grothe/LowerAnalytic/JointContractions.lean} supplying the direct
normalization and \code{Circular.joint\_sixth} API of the actual Gaussian
budget proof. The current cover has $8{,}756$ boxes, with $306$ spectral
and $45$ retained Stein nodes, including $41$ joint-moment nodes. Its
complete native Lean check passes at $B=0.7134218$, with strict endpoint
$K_G^{\mathbb C}>1.401695322458$. The downstream lower theorem retains
only \code{NodeCaps}, the previously documented Gaussian potential-node
bounds. Neither moment inequality alone proves the unrestricted
actual-response conjecture.

\subsection{Retaining the upper contraction bound and transporting a certificate}
\label{sec:joint-sixth-transport}
There is a further improvement to the two quoted moment bounds. Write
$s=\operatorname{tenS}(f)$, $h=\operatorname{tenH}(f)$,
$V=\operatorname{tenV}(f)$, $W=\operatorname{tenU2}(f)$ and
$D=\operatorname{tenDU}(f)$. Besides $D\le\sqrt{VW}$, one has
$0\le V\le s\sqrt s$ and $0\le W\le h\sqrt s$.
For every $\eta\ge0$, completion of the square with coefficient $20+\eta$
therefore gives
\[
 -20V+\frac{165}2D
 \le\eta V+\frac{27225}{16(20+\eta)}W
 \le\eta s\sqrt s+\frac{27225}{16(20+\eta)}h\sqrt s.
\]
The first inequality is precisely the nonnegativity of
\[
 (20+\eta)\left(\sqrt V-
 \frac{165}{4(20+\eta)}\sqrt W\right)^2
 +\frac{165}2(\sqrt{VW}-D).
\]
The same normalization, square-root bounds and Young inequality as above give,
for all $\eta\ge0$ and $\beta>0$,
\begin{equation}\label{eq:joint-sixth-eta}
 \E|q|^6\le6+9\chi+x\left\{
 \left(\frac{135\beta}2+2\eta\right)s_0+
 \left(198+\frac{27225}{8(20+\eta)}+
       \frac{135}{2\beta}\right)h_0\right\}.
\end{equation}
At $\eta=0$ this is \eqref{eq:joint-sixth-formal-young}.
Setting $\beta=\alpha+4/15$ and $\eta=165\alpha/4-9$ yields
\begin{equation}\label{eq:joint-sixth-shifted}
 \boxed{\quad
 \E|q|^6\le6+9\chi+x\left\{
 150\alpha s_0+\left(198+\frac{150}{\alpha+4/15}\right)h_0
 \right\},\qquad\alpha\ge\frac{12}{55}.
 \quad}
\end{equation}
It has the same $s_0$ coefficient as
\eqref{eq:joint-sixth-formal-triangle} and a strictly smaller $h_0$
coefficient. No concentration or coordinate-independence restriction is added.

\paragraph{An exact transport rule.}
Suppose the stronger tensor budget permits increasing the moment-cut weight
from $\zeta$ to $\zeta+d$, with $d\ge0$. If the old majorant has coefficients
$(a_0,a_1,a_2,a_4,a_6)$, replace them by
\[
 (a_0-12d,\ a_1,\ a_2,\ a_4+9d,\ a_6-d).
\]
The pointwise change in the majorant equals the pointwise change in the cut,
so the old whole-half-line majorant proof transports exactly. Provided
$d\le a_6$, the new sixth coefficient is still admissible. At normalized
scale its scalar moment price improves by exactly
\[
 J_{\rm new}(x)=J_{\rm old}(x)-dC_6x^3.
\]
This identity and the pointwise transport are proved in
\code{Grothe/Lower/MomentLift.lean}. The stronger budget must still be
accepted; the transport rule does not license an arbitrary increase of $d$.

For example, the two existing triangle nodes with $\alpha=1/10$ can use
\eqref{eq:joint-sixth-formal-young} at $\beta=1/8$ and increase their weight
by a factor $16/9$. The old bracket is $15s_0+1698h_0$, while the new
weighted bracket is $15s_0+(3229/2)h_0$. The freed coefficient is
$(167/2)h_0$, whose Gram representation is
\[
 h_0=\operatorname{tr}\!\left[
 \begin{pmatrix}2/3&8/3\\8/3&32/3\end{pmatrix}G_1\right]+8y.
\]
The displayed matrix is positive semidefinite of rank one. Adding its
nonnegative multiple to the old budget certificate proves acceptance of
the strengthened weight. Fresh exact replays of both transported nodes
and the complete $8{,}679$-box partition passed. Seventeen boxes improve,
by at most $4.091\cdot10^{-7}$ in this replay; other boxes remain binding,
so the displayed global lower endpoint of that package is unchanged.
A subsequent exact reselection and subdivision produces $8{,}756$ boxes with
$B=0.7134218$ and the strict external endpoint
$K_G^{\mathbb C}>1.401695322458$, under the original joint-moment package's
source inputs. Both modified nodes and the complete new cover pass the unchanged
exact checkers. The full joint-node, range and cover port now also passes
the native Lean check at this endpoint. The Gaussian potential-node bounds
remain the sole documented analytic input to the lower theorem.

The global strengthening \eqref{eq:joint-sixth-eta} and its shifted corollary
are formal theorems \code{JointSixth.gram\_j\_eta} and
\code{JointSixth.gram\_j2\_shifted}; \code{joint\_eta\_budget} connects the
former to the actual spectral budget. These introduce no additional Gaussian or
tensor assumptions. The lifted numerical package above uses the earlier Young
family; its new endpoint does not rely on the further eta strengthening.

\subsection{A homogeneous conic certificate for the moment cost}
\label{sec:moment-cost-cone}
The two parameters in \eqref{eq:joint-sixth-eta} need not be searched on an
outer grid. Let $\omega,u,v>0$ and choose real $a_s,a_h,r_0,r_1$ satisfying
\[
 \begin{pmatrix}u&135\omega/2\\135\omega/2&r_0\end{pmatrix}\succeq0,
 \qquad
 \begin{pmatrix}v&165\omega/2\\165\omega/2&r_1\end{pmatrix}\succeq0,
\]
\[
 v\ge40\omega,\qquad
 u+v\le a_s+40\omega,\qquad
 r_0+r_1\le a_h-198\omega.
\]
Then every normalized genuine circular cubic satisfies
\[
 \omega(\E|q|^6-6-9\chi)
 \le x(a_s s_0+a_h h_0).
\]
Indeed, take $\beta=2u/(135\omega)$ and
$\eta=v/(2\omega)-20$. The determinant tests give
\[
 \omega A_{\eta,\beta}=u+v-40\omega\le a_s,
 \qquad
 \omega B_{\eta,\beta}
 =198\omega+\frac{18225\omega^2}{4u}
                +\frac{27225\omega^2}{4v}
 \le198\omega+r_0+r_1\le a_h.
\]
For a node with $\omega=\zeta x_*$ and $x\le x_*$, this yields
\[
 \zeta(\E|q|^6-6-9\chi)\le a_s s_0+a_h h_0.
\]
Positivity of the multipliers and diagonal entries, and the range $x\le x_*$,
are explicit conditions; no division by a possibly zero quantity is used.
Nodes without a moment correction retain the separate zero-weight case.

This is a convex conic formulation of the moment-budget proposal: two
$2\times2$ PSD blocks and linear inequalities replace the outer
$(\eta,\beta)$ search. The reserved vector is linear in $(a_s,a_h)$, so it
can be optimized jointly with the contraction-Gram multipliers. It remains a
sufficient certificate family, not a characterization of attainable response laws.

\code{Grothe/LowerAnalytic/MomentCone.lean} proves the real cone theorem,
the node-range theorem, and soundness of an exact rational Boolean checker
(\code{check\_sound}, \code{gram\_node\_of\_check}). Two explicit rational
cone points, including the lifted nodes' $\beta=1/8$ point and the shifted
$\alpha=1/2$ point, are checked by kernel arithmetic. The latter has
$s_0$ coefficient $75$ and $h_0$ coefficient $9054/23$, compared with $498$
in the old triangle bound at the same $s_0$ coefficient. This coefficient
improvement alone is not an endpoint improvement.


\subsection{Exact optimization inside the moment-cost cone}
\label{sec:moment-cost-boundary}
The conic formulation itself has an exact two-branch scalar boundary. For a
fixed positive weight $w$ and positive $s_0$ cost $A$, the least admissible
$h_0$ cost is
\begin{equation}\label{eq:moment-cost-boundary}
 B_{\min}(w,A)=
 \begin{cases}
 \displaystyle\frac{11781}{32}w+\frac{18225w^2}{4A},&11A\le360w,\\[5pt]
 \displaystyle198w+\frac{22500w^2}{A+40w},&11A\ge360w.
 \end{cases}
\end{equation}
The formulas agree at $11A=360w$. This optimizes the sufficient certificate
cone, not the space of actual bounded-response laws.

Here is a direct proof retaining every constraint. Write the two positive
first diagonals as $u,v$, and the remaining diagonals as $r_0,r_1$.
Pairing the first PSD matrix against the second with its off-diagonal
negated gives
\[
 ur_1+vr_0\ge2(135w/2)(165w/2).
\]
Adding the two determinant bounds yields
\[
 (u+v)(r_0+r_1)\ge(150w)^2.
\]
Since $u+v\le A+40w$ and $r_0+r_1\le B-198w$, this proves the second
expression in \eqref{eq:moment-cost-boundary} as a universal lower bound.
It is attained when $11A\ge360w$ by
\[
 (u,v)=\frac{A+40w}{20}(9,11),\qquad
 (r_0,r_1)=\frac{w^2}{A+40w}(10125,12375).
\]
The threshold is precisely the condition $v\ge40w$ for this point.

For $11A\le360w$, put $\gamma=135w/(2A)\ge33/16$.
Testing the first PSD matrix on $(-\gamma,1)$ and the second on
$(-33/16,1)$, and using
\[
 (A-u)\gamma^2-(v-40w)(33/16)^2
 =\gamma^2(A+40w-u-v)
 +(\gamma^2-(33/16)^2)(v-40w)\ge0,
\]
gives
\[
 r_0+r_1\ge\frac{18225w^2}{4A}+\frac{5445}{32}w.
\]
After adding $198w$ this is the first branch. It is attained by
\[
 (u,v)=(A,40w),\qquad
 (r_0,r_1)=\left(\frac{18225w^2}{4A},\frac{5445w}{32}\right).
\]
The first branch is the original Young family; the second is the shifted
family, and both have rational attaining recipes for rational $w,A$.
Thus the strengthened two-parameter family has been completely optimized
inside this scalar cone. An outer parameter grid cannot produce a better
cost than \eqref{eq:moment-cost-boundary}.

The Lean proofs appear in
\code{Grothe/LowerAnalytic/MomentConeBoundary.lean}.
They establish necessity via \code{Feasible.boundary\_lower} and attainment
via \code{boundary\_attained}, together with the Schur and constrained-branch
inequalities. The proofs introduce no analytic hypotheses or native axioms.

The same module also proves \code{boundary\_moment} and \code{boundary\_node}: the optimal cost is applied directly to every actual normalized circular cubic, and to its moment correction whenever the checked kurtosis range holds.


In fact the auxiliary PSD variables can be eliminated completely.
\code{exists\_feasible\_iff} proves that, for $w,A>0$, existence of positive
$u,v$ and any $r_0,r_1$ satisfying the conic tests is equivalent to
$B\ge B_{\min}(w,A)$. The exact rational checker
\code{boundaryCheck} therefore needs only the three numbers $(w,A,B)$.
\code{boundaryChecked\_node} proves that its acceptance, together with
$w=\zeta x_*$, $\zeta,x_*>0$ and $x\le x_*$, implies
\[
 \zeta(\E|q|^6-6-9\chi)\le A s_0+B h_0
\]
for the actual circular cubic. Thus neither an outer parameter grid nor
stored auxiliary conic variables are needed for this moment-cost interface.
The remaining contraction reserve and pointwise node majorant still require
their own exact acceptance.
